\section{Тестирование}

Для проверки работоспособности предложенного подхода рассмотрим две характерные задачи встречи: задачу коррекции орбиты и задачу фазирования. Эти постановки отличаются структурой вектора конечного рассогласования $\Delta \mathbf{y}_f$. В задаче коррекции преобладают изменения формы и ориентации орбиты, тогда как в задаче фазирования основной вклад связан с изменением положения КА вдоль орбиты.

Во всех расчётах используется одна и та же модель пассивного движения. Она включает разложение гравитационного потенциала Земли до степени и порядка $32 \times 32$, сопротивление атмосферы, притяжение третьих тел и давление солнечного излучения. Параметры КА и ограничения на управление приведены в таблице \ref{tab::problem_correction_params}. Эффективные площади КА в модели атмосферного сопротивления и в модели давления солнечного излучения приняты одинаковыми. Для орбиты с большой полуосью $6900$ км максимальная продолжительность манёвра $\Delta t_{\text{max}} = 500$ с соответствует примерно $30^\circ$ дуги истинной долготы. Ограничение на характеристическую скорость одного манёвра определялось по формуле
\begin{equation}
    \Delta V_{\text{max}} = \frac{F}{m}\Delta t_{\text{max}},
\end{equation}
где $F$ --- модуль тяги, $m$ --- масса КА.

\begin{table}[!h]
\centering
\caption{Параметры КА и ограничения}
\label{tab::problem_correction_params}
\renewcommand{\arraystretch}{1.25}
\begin{tabular}{|c|c|c|c|c|c|c|}
\hline
$m$, кг & 
$F$, Н & 
$S_{\text{drag}} / m$, м$^2$/кг &
$S_{\text{srp}} / m$, м$^2$/кг &
$l_{\text{min}}$ &
$\Delta V_{\text{max}}$, м/c & 
$\Delta t_{\text{max}}$ \\
\hline
500 & 10 & 0.015 & 0.015 & 180 & 10 & 500 \\
\hline
\end{tabular}
\end{table}

\subsection{Задача коррекции орбиты}

Под задачей коррекции орбиты будем понимать частный случай задачи встречи, в котором основное рассогласование орбитальных элементов в конечный момент времени связано с формой и ориентацией орбиты, а изменение фазового положения КА не является определяющим фактором.

Параметры начальной орбиты, пассивного прогноза и целевой орбиты приведены в таблице \ref{tab::problem_correction_orb}. Для наглядности данные указаны в кеплеровых элементах. Продолжительность перелёта составляет $3$ дня. Целевая орбита задаётся как возмущение пассивного прогноза в конечный момент времени, поэтому в последней строке таблицы приведён соответствующий вектор конечного рассогласования $\Delta \mathbf{y}_f$.

\begin{table}[!h]
\centering
\caption{Параметры задачи коррекции орбиты}
\label{tab::problem_correction_orb}
\renewcommand{\arraystretch}{1.25}
\begin{tabular}{|c|c|c|c|c|c|c|}
\hline
 & $a$, км & $e$ & $\omega$, $^\circ$ & $i$, $^\circ$ & $\Omega$, $^\circ$ & $M$, $^\circ$ \\
\hline
$\mathbf{y}_i$ 
& 6900 & 0.001 & 0 & 53 & 0 & 0 \\
\hline
$\mathbf{y}_{\text{pass}}$
& 6899.7403 & 0.000376 & 210.4372 & 52.9657 & 346.2976 & 339.9065 \\
\hline
$\mathbf{y}_t$
& 6900.7403 & 0.010376 & 211.4372 & 53.0657 & 346.3976 & 340.9065 \\
\hline
$\Delta \mathbf{y}_f$
& 1 & 0.01 & 1 & 0.1 & 0.1 & 1 \\
\hline
\end{tabular}
\end{table}

Результаты решения задачи в импульсной постановке приведены в таблице \ref{tab::problem_correction_impulse_iterations}. Начальное приближение содержит шесть импульсов и уже обеспечивает выполнение конечных условий с невязкой порядка $10^{-2}$. Последующее нелинейное уточнение быстро уменьшает конечное рассогласование: после первой итерации его норма снижается до $1.7260 \cdot 10^{-4}$, а после четвёртой --- до $3.8395 \cdot 10^{-9}$. При этом суммарная характеристическая скорость практически не меняется и в окончательном импульсном решении составляет $42.7402$ м/с. Это показывает, что построенное начальное приближение не только близко к допустимому решению, но и хорошо согласовано с последующим нелинейным уточнением.

\renewcommand{\arraystretch}{1.25}
\begin{longtable}{|c|c|c|c|c|c|c|}
\caption{Импульсное решение задачи коррекции орбиты}
\label{tab::problem_correction_impulse_iterations}\\
\hline
{} & $j$ & $\Delta V_r$, м/с & $\Delta V_t$, м/с & $\Delta V_n$, м/с
& $\|\Delta \mathbf V\|$, м/с
& $L$, вит. \\
\hline
\endfirsthead

\hline
{} & $j$ & $\Delta V_r$, м/с & $\Delta V_t$, м/с & $\Delta V_n$, м/с
& $\|\Delta \mathbf V\|$, м/с
& $L$, вит. \\
\hline
\endhead

\multirow{7}{*}{\rotatebox[origin=c]{90}{\shortstack{Начальное\\приближение}}}
& 1 & -0.6315 & -3.1649 &  0.8318 & 3.3328 & 20.0223 \\*
& 2 & -0.0152 &  3.2767 & -1.8670 & 3.7713 & 21.6030 \\*
& 3 & -1.4940 & -7.8390 &  2.2932 & 8.3031 & 23.0486 \\*
& 4 & -1.4406 & -7.7933 &  2.6604 & 8.3599 & 24.0599 \\*
& 5 &  1.7602 &  8.1216 & -4.6034 & 9.5000 & 24.5930 \\*
& 6 &  2.6304 &  8.0747 & -4.2579 & 9.5000 & 25.5895 \\*
\cline{2-7}
& \multicolumn{6}{c|}{$\mathcal{J} = 42.7670$ м/с, \quad $\|\Delta \mathbf{y}\| = 9.4938 \cdot 10^{-3}$} \\
\hline

\multirow{7}{*}{\rotatebox[origin=c]{90}{Итерация 1}}
& 1 & -0.6167 & -2.7385 &  0.8597 & 2.9358 & 20.0223 \\*
& 2 &  0.0507 &  3.6760 & -1.8524 & 4.1167 & 21.6030 \\*
& 3 & -1.4921 & -8.0310 &  2.3067 & 8.4879 & 23.0486 \\*
& 4 & -1.4439 & -8.1647 &  2.6676 & 8.7099 & 24.0599 \\*
& 5 &  1.8219 &  7.9424 & -4.5930 & 9.3540 & 24.5930 \\*
& 6 &  2.6904 &  7.6649 & -4.2489 & 9.1674 & 25.5895 \\*
\cline{2-7}
& \multicolumn{6}{c|}{$\mathcal{J} = 42.7717$ м/с, \quad $\|\Delta \mathbf{y}\| = 1.7260 \cdot 10^{-4}$} \\
\hline

\multirow{7}{*}{\rotatebox[origin=c]{90}{Итерация 2}}
& 1 & -0.6163 & -2.7507 &  0.9211 & 2.9655 & 20.0223 \\*
& 2 &  0.0491 &  3.6629 & -1.8019 & 4.0824 & 21.6030 \\*
& 3 & -1.4913 & -8.0264 &  2.3316 & 8.4902 & 23.0486 \\*
& 4 & -1.4429 & -8.1560 &  2.6764 & 8.7044 & 24.0599 \\*
& 5 &  1.8205 &  7.9453 & -4.5555 & 9.3378 & 24.5930 \\*
& 6 &  2.6891 &  7.6749 & -4.2159 & 9.1602 & 25.5895 \\*
\cline{2-7}
& \multicolumn{6}{c|}{$\mathcal{J} = 42.7405$ м/с, \quad $\|\Delta \mathbf{y}\| = 5.3926 \cdot 10^{-6}$} \\
\hline

\multirow{7}{*}{\rotatebox[origin=c]{90}{Итерация 3}}
& 1 & -0.6163 & -2.7505 &  0.9222 & 2.9657 & 20.0223 \\*
& 2 &  0.0492 &  3.6632 & -1.8010 & 4.0822 & 21.6030 \\*
& 3 & -1.4913 & -8.0265 &  2.3320 & 8.4904 & 23.0486 \\*
& 4 & -1.4430 & -8.1562 &  2.6765 & 8.7046 & 24.0599 \\*
& 5 &  1.8206 &  7.9453 & -4.5547 & 9.3375 & 24.5930 \\*
& 6 &  2.6892 &  7.6748 & -4.2152 & 9.1598 & 25.5895 \\*
\cline{2-7}
& \multicolumn{6}{c|}{$\mathcal{J} = 42.7401$ м/с, \quad $\|\Delta \mathbf{y}\| = 1.1370 \cdot 10^{-7}$} \\
\hline

\multirow{7}{*}{\rotatebox[origin=c]{90}{Итерация 4}}
& 1 & -0.6163 & -2.7505 &  0.9221 & 2.9657 & 20.0223 \\*
& 2 &  0.0492 &  3.6632 & -1.8010 & 4.0822 & 21.6030 \\*
& 3 & -1.4913 & -8.0265 &  2.3320 & 8.4904 & 23.0486 \\*
& 4 & -1.4430 & -8.1562 &  2.6765 & 8.7046 & 24.0599 \\*
& 5 &  1.8206 &  7.9453 & -4.5548 & 9.3375 & 24.5930 \\*
& 6 &  2.6892 &  7.6748 & -4.2153 & 9.1598 & 25.5895 \\*
\cline{2-7}
& \multicolumn{6}{c|}{$\mathcal{J} = 42.7402$ м/с, \quad $\|\Delta \mathbf{y}\| = 3.8395 \cdot 10^{-9}$} \\
\hline

\end{longtable}

Решение той же задачи с учётом конечной длительности работы двигателя приведено в таблице \ref{tab::problem_correction_low_thrust_iterations}. В этом случае каждый импульс заменяется участком движения с постоянным управляющим ускорением в ОСК. Все найденные участки удовлетворяют ограничению на длительность манёвра: максимальная продолжительность включения ДУ в окончательном решении составляет $425.74$ с, что меньше заданного значения $\Delta t_{\text{max}}=500$ с.

\renewcommand{\arraystretch}{1.25}
\begin{longtable}{|c|c|c|c|c|c|c|c|c|}
\caption{Решение задачи коррекции орбиты с конечной тягой}
\label{tab::problem_correction_low_thrust_iterations}\\
\hline
{} & $j$
& $a_r$, м/с$^2$
& $a_t$, м/с$^2$
& $a_n$, м/с$^2$
& $\|\Delta \mathbf V\|$, м/с
& $L_{\mathrm{beg}}$, вит.
& $L_{\mathrm{end}}$, вит.
& $\Delta t$, с \\
\hline
\endfirsthead

\hline
{} & $j$
& $a_r$, м/с$^2$
& $a_t$, м/с$^2$
& $a_n$, м/с$^2$
& $\|\Delta \mathbf V\|$, м/с
& $L_{\mathrm{beg}}$, вит.
& $L_{\mathrm{end}}$, вит.
& $\Delta t$, с \\
\hline
\endhead

\multirow{7}{*}{\rotatebox[origin=c]{90}{\shortstack{Начальное\\приближение}}}
& 1 & -0.0033 & -0.0191 &  0.0050 & 4.9604 & 22.0149 & 22.0585 & 248.02 \\*
& 2 & -0.0036 & -0.0189 &  0.0055 & 7.3853 & 23.0160 & 23.0808 & 369.26 \\*
& 3 &  0.0021 &  0.0171 & -0.0101 & 5.7246 & 23.5706 & 23.6213 & 286.23 \\*
& 4 & -0.0034 & -0.0186 &  0.0064 & 7.4328 & 24.0274 & 24.0923 & 371.64 \\*
& 5 &  0.0037 &  0.0171 & -0.0097 & 8.4690 & 24.5553 & 24.6309 & 423.45 \\*
& 6 &  0.0055 &  0.0170 & -0.0090 & 8.4964 & 25.5516 & 25.6275 & 424.82 \\*
\cline{2-9}
& \multicolumn{8}{c|}{$\mathcal{J} = 42.4684$ м/с, \quad $\|\Delta \mathbf{y}\| = 1.2539 \cdot 10^{-2}$} \\
\hline

\multirow{7}{*}{\rotatebox[origin=c]{90}{Итерация 1}}
& 1 & -0.0032 & -0.0188 &  0.0059 & 4.3281 & 22.0177 & 22.0557 & 216.41 \\*
& 2 & -0.0033 & -0.0189 &  0.0055 & 7.4776 & 23.0156 & 23.0812 & 373.88 \\*
& 3 &  0.0024 &  0.0177 & -0.0090 & 6.5064 & 23.5671 & 23.6248 & 325.32 \\*
& 4 & -0.0029 & -0.0189 &  0.0058 & 8.2932 & 24.0237 & 24.0960 & 414.66 \\*
& 5 &  0.0041 &  0.0170 & -0.0097 & 8.5482 & 24.5549 & 24.6313 & 427.41 \\*
& 6 &  0.0065 &  0.0162 & -0.0098 & 7.8212 & 25.5546 & 25.6245 & 391.06 \\*
\cline{2-9}
& \multicolumn{8}{c|}{$\mathcal{J} = 42.9747$ м/с, \quad $\|\Delta \mathbf{y}\| = 2.2008 \cdot 10^{-4}$} \\
\hline

\multirow{7}{*}{\rotatebox[origin=c]{90}{Итерация 2}}
& 1 & -0.0032 & -0.0187 &  0.0062 & 4.3542 & 22.0176 & 22.0558 & 217.71 \\*
& 2 & -0.0033 & -0.0189 &  0.0057 & 7.4915 & 23.0156 & 23.0813 & 374.58 \\*
& 3 &  0.0025 &  0.0178 & -0.0088 & 6.4691 & 23.5673 & 23.6246 & 323.46 \\*
& 4 & -0.0029 & -0.0189 &  0.0058 & 8.2968 & 24.0237 & 24.0960 & 414.84 \\*
& 5 &  0.0041 &  0.0171 & -0.0096 & 8.5155 & 24.5551 & 24.6311 & 425.78 \\*
& 6 &  0.0065 &  0.0162 & -0.0097 & 7.7961 & 25.5547 & 25.6244 & 389.80 \\*
\cline{2-9}
& \multicolumn{8}{c|}{$\mathcal{J} = 42.9233$ м/с, \quad $\|\Delta \mathbf{y}\| = 6.7293 \cdot 10^{-6}$} \\
\hline

\multirow{7}{*}{\rotatebox[origin=c]{90}{Итерация 3}}
& 1 & -0.0032 & -0.0187 &  0.0062 & 4.3550 & 22.0176 & 22.0559 & 217.75 \\*
& 2 & -0.0033 & -0.0189 &  0.0057 & 7.4917 & 23.0156 & 23.0813 & 374.59 \\*
& 3 &  0.0025 &  0.0178 & -0.0088 & 6.4679 & 23.5673 & 23.6246 & 323.40 \\*
& 4 & -0.0029 & -0.0189 &  0.0058 & 8.2965 & 24.0237 & 24.0960 & 414.82 \\*
& 5 &  0.0041 &  0.0171 & -0.0096 & 8.5148 & 24.5551 & 24.6311 & 425.74 \\*
& 6 &  0.0065 &  0.0162 & -0.0097 & 7.7958 & 25.5547 & 25.6244 & 389.79 \\*
\cline{2-9}
& \multicolumn{8}{c|}{$\mathcal{J} = 42.9217$ м/с, \quad $\|\Delta \mathbf{y}\| = 1.0272 \cdot 10^{-8}$} \\
\hline

\multirow{7}{*}{\rotatebox[origin=c]{90}{Итерация 4}}
& 1 & -0.0032 & -0.0187 &  0.0062 & 4.3550 & 22.0176 & 22.0559 & 217.75 \\*
& 2 & -0.0033 & -0.0189 &  0.0057 & 7.4918 & 23.0156 & 23.0813 & 374.59 \\*
& 3 &  0.0025 &  0.0178 & -0.0088 & 6.4679 & 23.5673 & 23.6246 & 323.39 \\*
& 4 & -0.0029 & -0.0189 &  0.0058 & 8.2964 & 24.0237 & 24.0960 & 414.82 \\*
& 5 &  0.0041 &  0.0171 & -0.0096 & 8.5147 & 24.5551 & 24.6311 & 425.74 \\*
& 6 &  0.0065 &  0.0162 & -0.0097 & 7.7958 & 25.5547 & 25.6244 & 389.79 \\*
\cline{2-9}
& \multicolumn{8}{c|}{$\mathcal{J} = 42.9216$ м/с, \quad $\|\Delta \mathbf{y}\| = 1.1431 \cdot 10^{-9}$} \\
\hline

\end{longtable}

Сопоставление таблиц \ref{tab::problem_correction_impulse_iterations} и \ref{tab::problem_correction_low_thrust_iterations} показывает, что при переходе к модели конечной тяги сохраняется общая структура найденного решения. В обоих случаях используется шесть манёвров, компактно сосредоточенных с $20$ по $25$ витки перелёта. Протяжённые участки работы ДУ располагаются в той же области по истинной долготе, что и импульсы исходного решения, а их длительности не превышают заданного ограничения $\Delta t_{\text{max}}=500$ с. Изменение суммарной характеристической скорости остаётся малым: она возрастает с $42.7402$ м/с до $42.9216$ м/с, то есть примерно на $0.42\%$. Конечное рассогласование после четырёх итераций составляет $1.1431 \cdot 10^{-9}$, что подтверждает сходимость процедуры уточнения и применимость построенного импульсного приближения для перехода к манёврам конечной длительности.
 
\subsection{Задача фазирования}

Под задачей фазирования будем понимать частный случай задачи встречи, в котором основная цель состоит в изменении положения КА вдоль орбиты при сравнительно малых изменениях остальных орбитальных элементов. Параметры задачи приведены в таблице \ref{tab::problem_phasing}. Продолжительность перелёта составляет $10$ дней. В отличие от предыдущего примера, здесь вектор конечного рассогласования имеет выраженную фазовую компоненту: изменение средней аномалии составляет $50^\circ$, тогда как изменения остальных элементов существенно меньше.

\begin{table}[!h]
\centering
\caption{Параметры задачи фазирования}
\label{tab::problem_phasing}
\renewcommand{\arraystretch}{1.25}
\begin{tabular}{|c|c|c|c|c|c|c|}
\hline
 & $a$, км & $e$ & $\omega$, $^\circ$ & $i$, $^\circ$ & $\Omega$, $^\circ$ & $M$, $^\circ$ \\
\hline
$\mathbf{y}_i$ 
& 6900 & 0.001 & 0 & 53 & 0 & 0 \\
\hline
$\mathbf{y}_{\text{pass}}$
& 6887.9941 & 0.001135 & 54.6669 & 52.8687 & 314.2723 & 219.9063 \\
\hline
$\mathbf{y}_t$
& 6888.0941 & 0.002135 & 54.7669 & 52.8787 & 314.2823 & 269.9063 \\
\hline
$\Delta \mathbf{y}_f$
& 0.1 & 0.001 & 0.1 & 0.01 & 0.01 & 50 \\
\hline
\end{tabular}
\end{table}

Результаты решения задачи фазирования в импульсной постановке приведены в таблице \ref{tab::problem_phasing_impulse_iterations}. Начальное приближение содержит четыре импульса и имеет суммарную характеристическую скорость $20.3883$ м/с. В ходе нелинейного уточнения происходит существенное перераспределение импульсов скорости, после чего значение функционала уменьшается до $16.9018$ м/с. Конечное рассогласование после четвёртой итерации составляет $7.9710 \cdot 10^{-9}$.

\renewcommand{\arraystretch}{1.25}
\begin{longtable}{|c|c|c|c|c|c|c|}
\caption{Импульсное решение задачи фазирования}
\label{tab::problem_phasing_impulse_iterations}\\
\hline
{} & $j$ & $\Delta V_r$, м/с & $\Delta V_t$, м/с & $\Delta V_n$, м/с
& $\|\Delta \mathbf V\|$, м/с
& $L$, вит. \\
\hline
\endfirsthead

\hline
{} & $j$ & $\Delta V_r$, м/с & $\Delta V_t$, м/с & $\Delta V_n$, м/с
& $\|\Delta \mathbf V\|$, м/с
& $L$, вит. \\
\hline
\endhead

\multirow{5}{*}{\rotatebox[origin=c]{90}{\shortstack{Начальное\\приближение}}}
& 1 & -0.5457 & -1.0567 & -2.0271 & 2.3502 & 86.6630 \\*
& 2 & -2.1820 & -4.4729 & -8.0921 & 9.5000 & 87.6586 \\*
& 3 & 0.5234 & 0.2070 & 1.4733 & 1.5772 & 147.1903 \\*
& 4 & 2.0658 & 0.8992 & 6.5862 & 6.9609 & 150.1875 \\*
\cline{2-7}
& \multicolumn{6}{c|}{$\mathcal{J} = 20.3883$ м/с, \quad $\|\Delta \mathbf{y}\| = 1.6957 \cdot 10^{-3}$} \\
\hline

\multirow{5}{*}{\rotatebox[origin=c]{90}{Итерация 1}}
& 1 & -0.7473 & -1.1005 & -0.4157 & 1.3937 & 86.6630 \\*
& 2 & -2.4236 & -4.4335 & -6.4309 & 8.1785 & 87.6586 \\*
& 3 & 0.2974 & 0.0750 & 0.9099 & 0.9602 & 147.1903 \\*
& 4 & 1.8597 & 1.0621 & 6.0094 & 6.3796 & 150.1875 \\*
\cline{2-7}
& \multicolumn{6}{c|}{$\mathcal{J} = 16.9119$ м/с, \quad $\|\Delta \mathbf{y}\| = 4.8929 \cdot 10^{-5}$} \\
\hline

\multirow{5}{*}{\rotatebox[origin=c]{90}{Итерация 2}}
& 1 & -0.7239 & -1.0966 & -0.4077 & 1.3758 & 86.6630 \\*
& 2 & -2.3960 & -4.4385 & -6.4224 & 8.1663 & 87.6586 \\*
& 3 & 0.3200 & 0.0887 & 0.9145 & 0.9729 & 147.1903 \\*
& 4 & 1.8802 & 1.0446 & 6.0140 & 6.3871 & 150.1875 \\*
\cline{2-7}
& \multicolumn{6}{c|}{$\mathcal{J} = 16.9021$ м/с, \quad $\|\Delta \mathbf{y}\| = 2.4993 \cdot 10^{-6}$} \\
\hline

\multirow{5}{*}{\rotatebox[origin=c]{90}{Итерация 3}}
& 1 & -0.7240 & -1.0966 & -0.4075 & 1.3757 & 86.6630 \\*
& 2 & -2.3961 & -4.4385 & -6.4221 & 8.1661 & 87.6586 \\*
& 3 & 0.3199 & 0.0886 & 0.9145 & 0.9729 & 147.1903 \\*
& 4 & 1.8802 & 1.0447 & 6.0140 & 6.3871 & 150.1875 \\*
\cline{2-7}
& \multicolumn{6}{c|}{$\mathcal{J} = 16.9018$ м/с, \quad $\|\Delta \mathbf{y}\| = 1.1866 \cdot 10^{-8}$} \\
\hline

\multirow{5}{*}{\rotatebox[origin=c]{90}{Итерация 4}}
& 1 & -0.7240 & -1.0966 & -0.4075 & 1.3757 & 86.6630 \\*
& 2 & -2.3961 & -4.4385 & -6.4222 & 8.1661 & 87.6586 \\*
& 3 & 0.3199 & 0.0886 & 0.9145 & 0.9729 & 147.1903 \\*
& 4 & 1.8802 & 1.0447 & 6.0140 & 6.3871 & 150.1875 \\*
\cline{2-7}
& \multicolumn{6}{c|}{$\mathcal{J} = 16.9018$ м/с, \quad $\|\Delta \mathbf{y}\| = 7.9710 \cdot 10^{-9}$} \\
\hline

\end{longtable}

Решение задачи фазирования с конечной тягой приведено в таблице \ref{tab::problem_phasing_low_thrust_iterations}. Как и в импульсной постановке, окончательное решение содержит четыре манёвра. При этом центры участков работы двигателя практически совпадают с истинными долготами, найденными в импульсной модели. Такое совпадение показывает, что для данной задачи импульсное решение хорошо определяет не только число манёвров, но и их расположение на интервале перелёта.

\renewcommand{\arraystretch}{1.25}
\begin{longtable}{|c|c|c|c|c|c|c|c|c|}
\caption{Решение задачи задачи фазирования с конечной тягой}
\label{tab::problem_phasing_low_thrust_iterations}\\
\hline
{} & $j$
& $a_r$, м/с$^2$
& $a_t$, м/с$^2$
& $a_n$, м/с$^2$
& $\|\Delta \mathbf V\|$, м/с
& $L_{\mathrm{beg}}$, вит.
& $L_{\mathrm{end}}$, вит.
& $\Delta t$, с \\
\hline
\endfirsthead

\hline
{} & $j$
& $a_r$, м/с$^2$
& $a_t$, м/с$^2$
& $a_n$, м/с$^2$
& $\|\Delta \mathbf V\|$, м/с
& $L_{\mathrm{beg}}$, вит.
& $L_{\mathrm{end}}$, вит.
& $\Delta t$, с \\
\hline
\endhead

\multirow{5}{*}{\rotatebox[origin=c]{90}{\shortstack{Начальное\\приближение}}}
& 1 & -0.0046 & -0.0090 & -0.0173 & 3.3924 & 86.6482 & 86.6780 & 169.62 \\*
& 2 & -0.0046 & -0.0094 & -0.0170 & 8.4962 & 87.6213 & 87.6958 & 424.81 \\*
& 3 & 0.0066 & 0.0026 & 0.0187 & 2.2768 & 147.1804 & 147.2004 & 113.84 \\*
& 4 & 0.0059 & 0.0026 & 0.0189 & 6.2325 & 150.1601 & 150.2151 & 311.63 \\*
\cline{2-9}
& \multicolumn{8}{c|}{$\mathcal{J} = 20.3979$ м/с, \quad $\|\Delta \mathbf{y}\| = 1.8480 \cdot 10^{-3}$} \\
\hline

\multirow{5}{*}{\rotatebox[origin=c]{90}{Итерация 1}}
& 1 & -0.0080 & -0.0138 & -0.0120 & 2.2250 & 86.6533 & 86.6728 & 111.25 \\*
& 2 & -0.0058 & -0.0111 & -0.0156 & 7.1909 & 87.6270 & 87.6901 & 359.54 \\*
& 3 & 0.0077 & 0.0028 & 0.0182 & 1.7218 & 147.1828 & 147.1980 & 86.09 \\*
& 4 & 0.0062 & 0.0031 & 0.0187 & 5.6863 & 150.1625 & 150.2127 & 284.32 \\*
\cline{2-9}
& \multicolumn{8}{c|}{$\mathcal{J} = 16.8240$ м/с, \quad $\|\Delta \mathbf{y}\| = 2.6790 \cdot 10^{-4}$} \\
\hline

\multirow{5}{*}{\rotatebox[origin=c]{90}{Итерация 2}}
& 1 & -0.0081 & -0.0139 & -0.0119 & 2.2207 & 86.6533 & 86.6728 & 111.04 \\*
& 2 & -0.0058 & -0.0111 & -0.0156 & 7.1809 & 87.6271 & 87.6900 & 359.04 \\*
& 3 & 0.0075 & 0.0027 & 0.0183 & 1.7169 & 147.1828 & 147.1980 & 85.84 \\*
& 4 & 0.0062 & 0.0032 & 0.0188 & 5.6854 & 150.1625 & 150.2127 & 284.27 \\*
\cline{2-9}
& \multicolumn{8}{c|}{$\mathcal{J} = 16.8039$ м/с, \quad $\|\Delta \mathbf{y}\| = 8.0427 \cdot 10^{-6}$} \\
\hline

\multirow{5}{*}{\rotatebox[origin=c]{90}{Итерация 3}}
& 1 & -0.0081 & -0.0139 & -0.0119 & 2.2203 & 86.6533 & 86.6728 & 111.02 \\*
& 2 & -0.0058 & -0.0111 & -0.0156 & 7.1804 & 87.6271 & 87.6900 & 359.02 \\*
& 3 & 0.0075 & 0.0027 & 0.0183 & 1.7170 & 147.1828 & 147.1980 & 85.85 \\*
& 4 & 0.0062 & 0.0032 & 0.0188 & 5.6854 & 150.1625 & 150.2127 & 284.27 \\*
\cline{2-9}
& \multicolumn{8}{c|}{$\mathcal{J} = 16.8031$ м/с, \quad $\|\Delta \mathbf{y}\| = 1.1913 \cdot 10^{-7}$} \\
\hline

\multirow{5}{*}{\rotatebox[origin=c]{90}{Итерация 4}}
& 1 & -0.0081 & -0.0139 & -0.0119 & 2.2203 & 86.6533 & 86.6728 & 111.02 \\*
& 2 & -0.0058 & -0.0111 & -0.0156 & 7.1804 & 87.6271 & 87.6900 & 359.02 \\*
& 3 & 0.0075 & 0.0027 & 0.0183 & 1.7170 & 147.1828 & 147.1980 & 85.85 \\*
& 4 & 0.0062 & 0.0032 & 0.0188 & 5.6854 & 150.1625 & 150.2127 & 284.27 \\*
\cline{2-9}
& \multicolumn{8}{c|}{$\mathcal{J} = 16.8031$ м/с, \quad $\|\Delta \mathbf{y}\| = 7.7198 \cdot 10^{-8}$} \\
\hline

\multirow{5}{*}{\rotatebox[origin=c]{90}{Итерация 5}}
& 1 & -0.0081 & -0.0139 & -0.0119 & 2.2203 & 86.6533 & 86.6728 & 111.02 \\*
& 2 & -0.0058 & -0.0111 & -0.0156 & 7.1804 & 87.6271 & 87.6900 & 359.02 \\*
& 3 & 0.0075 & 0.0027 & 0.0183 & 1.7170 & 147.1828 & 147.1980 & 85.85 \\*
& 4 & 0.0062 & 0.0032 & 0.0188 & 5.6854 & 150.1625 & 150.2127 & 284.27 \\*
\cline{2-9}
& \multicolumn{8}{c|}{$\mathcal{J} = 16.8031$ м/с, \quad $\|\Delta \mathbf{y}\| = 7.0310 \cdot 10^{-8}$} \\
\hline

\multirow{5}{*}{\rotatebox[origin=c]{90}{Итерация 6}}
& 1 & -0.0081 & -0.0139 & -0.0119 & 2.2203 & 86.6533 & 86.6728 & 111.02 \\*
& 2 & -0.0058 & -0.0111 & -0.0156 & 7.1804 & 87.6271 & 87.6900 & 359.02 \\*
& 3 & 0.0075 & 0.0027 & 0.0183 & 1.7170 & 147.1828 & 147.1980 & 85.85 \\*
& 4 & 0.0062 & 0.0032 & 0.0188 & 5.6854 & 150.1625 & 150.2127 & 284.27 \\*
\cline{2-9}
& \multicolumn{8}{c|}{$\mathcal{J} = 16.8031$ м/с, \quad $\|\Delta \mathbf{y}\| = 8.5511 \cdot 10^{-9}$} \\
\hline

\end{longtable}

Суммарная характеристическая скорость в решении с конечной тягой составляет $16.8031$ м/с, тогда как в импульсной постановке она равна $16.9018$ м/с. Различие между этими значениями не превышает $1\%$, что указывает на близость двух решений. Небольшое уменьшение $\mathcal{J}$ при переходе к конечной тяге связано с тем, что манёвр конечной длительнности не является точной заменой мгновенного импульса: управление действует на конечном интервале времени, и его вклад в конечное состояние формируется вдоль всего активного участка.

Следует отметить, что для решения с конечной тягой потребовалось больше итераций, чем в импульсном случае. Это связано с большей нелинейностью задачи фазирования: малые изменения большой полуоси и среднего движения на интервале в $10$ дней приводят к накоплению фазового отклонения, поэтому конечное условие оказывается более чувствительным к параметрам протяжённых манёвров. Тем не менее алгоритм сходится устойчиво: норма конечного рассогласования уменьшается от $1.8480 \cdot 10^{-3}$ для начального приближения до $8.5511 \cdot 10^{-9}$ в окончательном решении. Все длительности включения ДУ остаются меньше $\Delta t_{\text{max}}=500$ с, что подтверждает выполнение заданных ограничений.

В совокупности оба примера показывают, что предложенная схема построения начального приближения и последующего нелинейного уточнения применима как к задаче коррекции орбиты, так и к задаче фазирования. В обоих случаях импульсное решение даёт информативную структуру манёвров, которая может быть использована при переходе к модели конечной тяги. При этом конечное рассогласование уменьшается до малых значений за несколько итераций, а изменение суммарной характеристической скорости при учёте конечной длительности манёвров остаётся небольшим.

\subsection{Пока хз}

